\documentclass[11pt]{article}
\usepackage[a4paper,margin=2.5cm]{geometry}
\setlength{\parskip}{0.6em}
\setlength{\parindent}{0pt}

\title{\textbf{Technology Statement}\\[4pt]
\large Clothing Classification on Fashion-MNIST}
\author{
Thijs Niekerk (SNR: 2125613) \\
Pjotr Martens (SNR: 2080478) \\
\\
Tilburg School of Economics and Management \\
Tilburg University}
\date{\today}

\begin{document}
\maketitle

This statement accompanies our submission for the Machine Learning Analytics group assignment on Fashion-MNIST classification, in line with the course's and Tilburg University's policy on the use of generative AI and other technology in coursework.

\textbf{Software and libraries.} All code was written in Python 3, using \texttt{numpy}, \texttt{matplotlib} and \texttt{scikit-learn} (Pedregosa et al., 2011). The Fashion-MNIST data and its loading format follow the description in the dataset's repository (Xiao, Rasul \& Vollgraf, 2017).

\textbf{Use of generative AI.} We used Claude (Anthropic), a generative AI assistant, during this project. Specifically, it was used to help draft the data-loading and model-training code in \texttt{code/}, to help typeset this report and the accompanying figures in \LaTeX{}, and to help draft an initial version of the report text based on the results our code produced. Every reported number and figure comes from actually running the submitted code; we reviewed, checked and edited the resulting code, analysis and text ourselves, and we take full responsibility for the correctness and originality of the final submission.

\textbf{Other sources.} Where we relied on external material beyond standard library documentation (for example the published Fashion-MNIST benchmark results used in Part 2e), this is cited in the report's reference list.

\vspace{1em}
We declare that this statement accurately reflects how this submission was produced.

\vspace{2em}
Signed: \underline{\hspace{4cm}} \hspace{2cm} \underline{\hspace{4cm}}

\end{document}
